\documentclass[11pt,a4paper]{article}
\usepackage[utf8]{inputenc}
\usepackage[T1]{fontenc}
\usepackage[portuguese]{babel}
\usepackage{amsmath,amssymb}
\usepackage{booktabs}
\usepackage{graphicx}
\usepackage[colorlinks=true,linkcolor=blue,citecolor=blue,urlcolor=blue]{hyperref}
\usepackage{geometry}
\geometry{margin=2.5cm}
\usepackage{tikz}
\usetikzlibrary{shapes,arrows.meta,positioning,fit,backgrounds}
\usepackage{tcolorbox}
\usepackage{enumitem}

\title{Protocolo de Treinamento e Metodologia Experimental — NODE v6}
\author{}
\date{}

\begin{document}

\maketitle

\begin{center}
\Large\textbf{Estudo Comparativo de Sinais de Excitação para Identificação de Sistemas via Neural ODEs}
\end{center}

\vspace{0.5cm}

\section{Introdução e Objetivo}

O presente estudo investiga a influência do tipo de sinal de excitação utilizado no treinamento sobre a capacidade de generalização de modelos baseados em \textbf{Neural Ordinary Differential Equations} (Neural ODEs) aplicados à identificação de um aeropêndulo.

A pergunta central é:

\begin{tcolorbox}
\textit{Dado um mesmo modelo e os mesmos hiperparâmetros, qual tipo de sinal de excitação no conjunto de treino produz o modelo NODE mais generalizável?}
\end{tcolorbox}

Para responder, seis experimentos (A--F) são conduzidos sob condições idênticas, variando-se exclusivamente o conjunto de dados de treinamento. A avaliação é feita em \textbf{simulação livre} (\textit{free-run}) sobre datasets de teste de sessões experimentais completamente distintas (\textit{cross-session}), cobrindo todos os cinco tipos de excitação.

\section{Descrição da Planta}

O sistema sob estudo é um \textbf{aeropêndulo de 1 grau de liberdade (1-DOF)}: uma haste articulada em torno de um eixo horizontal, acionada por um motor com hélice em uma extremidade e balanceada por um contrapeso na outra.

\subsection{Parâmetros Físicos Conhecidos}

\begin{table}[hbt!]
\centering
\caption{Parâmetros Físicos Conhecidos}
\vspace{0.2cm}
\begin{tabular}{lccc}
\toprule
\textbf{Parâmetro} & \textbf{Símbolo} & \textbf{Valor} & \textbf{Unidade} \\
\midrule
Massa da hélice + motor & $m_1$ & 0.122 & kg \\
Comprimento ao motor & $L_1$ & 0.39 & m \\
Massa do contrapeso & $m_2$ & 0.055 & kg \\
Comprimento ao contrapeso & $L_2$ & 0.347 & m \\
Aceleração gravitacional & $g$ & 9.81 & m/s$^2$ \\
\bottomrule
\end{tabular}
\end{table}

\subsection{Variáveis do Sistema}

\begin{table}[hbt!]
\centering
\caption{Variáveis do Sistema}
\vspace{0.2cm}
\begin{tabular}{lcl}
\toprule
\textbf{Variável} & \textbf{Símbolo} & \textbf{Descrição} \\
\midrule
Entrada & $u(t)$ & Comando do motor (\% do PWM), normalizado para $[-1, 1]$ \\
Saída & $\theta(t)$ & Ângulo da haste (rad) \\
Estado & $\mathbf{x}(t) = [\theta,\, \dot{\theta}]^\top$ & Posição angular e velocidade angular \\
\bottomrule
\end{tabular}
\end{table}

\subsection{Instrumentação e Aquisição}

\begin{itemize}
    \item \textbf{Sensor de ângulo}: encoder ou potenciômetro no eixo de rotação
    \item \textbf{Taxa de amostragem nativa}: $f_s = 100$ Hz ($\Delta t_{\text{raw}} = 0{,}010$ s)
    \item \textbf{Taxa efetiva após decimação}: $f_s' = 50$ Hz ($\Delta t = 0{,}020$ s)
    \item \textbf{Controlador}: microcontrolador (Arduino/STM32) com comunicação serial
\end{itemize}

\section{Modelo — Physics-Informed Neural ODE Assimétrico}

\subsection{Formulação Dinâmica}

O modelo segue a forma de uma equação diferencial ordinária parametrizada, onde os parâmetros físicos desconhecidos são estimados via otimização baseada em gradiente:
%
\begin{equation}
\dot{\mathbf{x}}(t) = f_\phi(\mathbf{x}(t),\, u(t))
\end{equation}

Especificamente, a dinâmica é descrita por:
%
\begin{align}
\dot{\theta} &= \dot{\theta} \\
\ddot{\theta} &= \frac{1}{J} \left[ \tau_{\text{motor}}(u) - \tau_{\text{grav}}(\theta) - \tau_{\text{atrito}}(\dot{\theta}) \right]
\end{align}

onde:

\begin{table}[hbt!]
\centering
\caption{Componentes de Torque}
\vspace{0.2cm}
\begin{tabular}{lll}
\toprule
\textbf{Torque} & \textbf{Expressão} & \textbf{Descrição} \\
\midrule
Motor & $\tau_{\text{motor}} = G_u(u) \cdot u \cdot |u|$ & Torque aerodinâmico (quadrático em $u$) \\
Gravitacional & $\tau_{\text{grav}} = (m_1 L_1 - m_2 L_2)\, g\, \sin\theta$ & Torque restaurador/desestabilizador \\
Atrito & $\tau_{\text{atrito}} = b(\dot{\theta}) \cdot \dot{\theta}$ & Atrito viscoso assimétrico \\
\bottomrule
\end{tabular}
\end{table}

\subsection{Assimetria Direção-Dependente}

Para capturar a assimetria física do sistema (diferenças no atrito e no empuxo entre rotação horária/anti-horária e entre comando positivo/negativo), os coeficientes de atrito e ganho do motor são modelados como funções suaves da direção:
%
\begin{align}
b(\dot{\theta}) &= \sigma(50\,\dot{\theta}) \cdot b^+ + \left[1 - \sigma(50\,\dot{\theta})\right] \cdot b^- \\
G_u(u) &= \sigma(50\,u) \cdot G_u^+ + \left[1 - \sigma(50\,u)\right] \cdot G_u^-
\end{align}

onde $\sigma(\cdot)$ é a função sigmóide logística e o fator 50 produz uma transição rápida (quase-degrau) próxima à origem.

\subsection{Parâmetros Treináveis}

Todos os parâmetros são parametrizados em escala logarítmica para garantir positividade:

\begin{table}[hbt!]
\centering
\caption{Parâmetros Treináveis}
\vspace{0.2cm}
\begin{tabular}{lcc}
\toprule
\textbf{Parâmetro} & \textbf{Notação Interna} & \textbf{Valor Inicial} \\
\midrule
Momento de inércia & $J = e^{\log J}$ & 1.0 \\
Atrito viscoso (positivo) & $b^+ = e^{\log b^+}$ & $e^{-1} \approx 0.368$ \\
Atrito viscoso (negativo) & $b^- = e^{\log b^-}$ & $e^{-1} \approx 0.368$ \\
Ganho do motor (positivo) & $G_u^+ = e^{\log G_u^+}$ & 1.0 \\
Ganho do motor (negativo) & $G_u^- = e^{\log G_u^-}$ & 1.0 \\
\bottomrule
\end{tabular}
\end{table}

\textbf{Total: 5 parâmetros treináveis.}

\subsection{Integração Numérica}

A integração da ODE é realizada pelo pacote \texttt{torchdiffeq} utilizando o método \textbf{Runge-Kutta de 4ª ordem (RK4)} a passo fixo, com interpolação linear do sinal de entrada $u(t)$ via \texttt{torch.searchsorted}.

\section{Dados Experimentais}

\subsection{Campanhas de Coleta (Rodadas)}

Os dados experimentais foram coletados em cinco sessões (rodadas) independentes, em datas distintas, cada uma potencialmente sob condições ambientais e de setup ligeiramente diferentes:

\begin{table}[hbt!]
\centering
\caption{Campanhas de Coleta (Rodadas)}
\vspace{0.2cm}
\begin{tabular}{clll}
\toprule
\textbf{Rodada} & \textbf{Data} & \textbf{Conteúdo Principal} & \textbf{Uso no v6} \\
\midrule
1 & 31/07 e 03/08 & Multi-seno, swept-sine, chirp (piloto) & --- (não utilizada) \\
2 & 04/08 & Multi-seno, seq-degraus, chirp & \textbf{Teste} \\
3 & 07/08 & Chirp, seq-degraus, multi-seno, curva & \textbf{Val} + \textbf{Treino (Exp E)} \\
4 & 19/08 & APRBS, multi-seno, swept-sine, degraus & \textbf{Teste} \\
5 (raiz) & 27/08 & APRBS, multi-seno, swept-sine, seq-degraus & \textbf{Treino (Exp A--D)} \\
\bottomrule
\end{tabular}
\end{table}

\subsection{Tipos de Sinais de Excitação}

\begin{table}[hbt!]
\centering
\caption{Tipos de Sinais de Excitação}
\vspace{0.2cm}
\begin{tabular}{lp{2cm}p{9.5cm}}
\toprule
\textbf{Tipo} & \textbf{Abreviatura} & \textbf{Característica Espectral} \\
\midrule
APRBS & APRBS & Sinal binário pseudo-aleatório de amplitude, espectro largo e rico em baixas frequências \\
Multi-seno & MultiSeno & Soma de senóides com frequências escolhidas, espectro discreto controlado \\
Swept-sine & SweptSine & Varredura contínua de frequência (chirp linear), cobertura espectral progressiva \\
Seq. degraus & Degraus & Degraus de amplitude variada, rico em conteúdo DC e transientes \\
Chirp & Chirp & Varredura logarítmica de frequência, ênfase proporcional em cada década \\
\bottomrule
\end{tabular}
\end{table}

\subsection{Inventário Completo dos Arquivos Utilizados}

\subsubsection*{Treino — Experimentos A--D (Rodada 5, 27/08)}

\begin{table}[hbt!]
\centering
\caption{Treino --- Experimentos A--D}
\vspace{0.2cm}
\begin{tabular}{llc}
\toprule
\textbf{Experimento} & \textbf{Arquivo} & \textbf{Tipo} \\
\midrule
A\_APRBS & \texttt{aprbs-1\_0827\_17-19.csv} & APRBS \\
 & \texttt{aprbs-2\_0827\_17-25.csv} & APRBS \\
 & \texttt{aprbs-3\_0827\_17-28.csv} & APRBS \\
 & \texttt{aprbs-4\_0827\_17-31.csv} & APRBS \\
\midrule
B\_MultiSeno & \texttt{multi-seno-1\_0827\_17-34.csv} & Multi-seno \\
 & \texttt{multi-seno-2\_0827\_17-37.csv} & Multi-seno \\
 & \texttt{multi-seno-3\_0827\_17-40.csv} & Multi-seno \\
 & \texttt{multi-seno-4\_0827\_17-43.csv} & Multi-seno \\
\midrule
C\_SweptSine & \texttt{swept-sine-1\_0827\_17-58.csv} & Swept-sine \\
 & \texttt{swept-sine-4\_0827\_18-00.csv} & Swept-sine \\
\midrule
D\_Degraus & \texttt{seq-degraus-1\_0827\_17-46.csv} & Seq-degraus \\
 & \texttt{seq-degraus-2\_0827\_17-49.csv} & Seq-degraus \\
 & \texttt{seq-degraus-3\_0827\_17-52.csv} & Seq-degraus \\
 & \texttt{seq-degraus-4\_0827\_17-55.csv} & Seq-degraus \\
\bottomrule
\end{tabular}
\end{table}

\subsubsection*{Treino — Experimento E (Rodada 3, 07/08)}

\begin{table}[hbt!]
\centering
\caption{Treino --- Experimento E}
\vspace{0.2cm}
\begin{tabular}{llc}
\toprule
\textbf{Experimento} & \textbf{Arquivo} & \textbf{Tipo} \\
\midrule
E\_Chirp & \texttt{RODADA-3/chirp-1\_0807\_16-34.csv} & Chirp \\
 & \texttt{RODADA-3/chirp-2\_0807\_17-07.csv} & Chirp \\
 & \texttt{RODADA-3/chirp-2\_0807\_17-09.csv} & Chirp \\
\bottomrule
\end{tabular}
\end{table}

\begin{tcolorbox}[colback=blue!5!white,colframe=blue!75!black,title=Nota]
Não há arquivos de chirp na Rodada 5 (raiz). Os chirps do Exp E são da Rodada 3, porém nenhum deles coincide com o arquivo de validação (\texttt{chirp-1\_0807\_16-32.csv}) nem com o de teste (\texttt{RODADA-2/chirp-1\_0804\_19-19.csv}).
\end{tcolorbox}

\subsubsection*{Treino — Experimento F: Mix (Controle)}

Concatenação sem duplicatas de todos os arquivos dos Experimentos A--E (17 arquivos no total).

\subsubsection*{Validação (\textit{early stopping})}

\begin{table}[hbt!]
\centering
\caption{Validação (\textit{early stopping})}
\vspace{0.2cm}
\begin{tabular}{lcc}
\toprule
\textbf{Arquivo} & \textbf{Rodada} & \textbf{Tipo} \\
\midrule
\texttt{RODADA-3/chirp-1\_0807\_16-32.csv} & 3 & Chirp \\
\bottomrule
\end{tabular}
\end{table}

Função: monitorar o RMSE em simulação livre a cada 50 épocas. O checkpoint com menor RMSE de validação é restaurado ao final do treinamento.

\subsubsection*{Teste (avaliação final, \textit{locked})}

\begin{table}[hbt!]
\centering
\caption{Teste (avaliação final, \textit{locked})}
\vspace{0.2cm}
\begin{tabular}{llc}
\toprule
\textbf{Tipo de Teste} & \textbf{Arquivo} & \textbf{Rodada} \\
\midrule
APRBS & \texttt{RODADA-4/aprbs-2\_0819\_18-51.csv} & 4 \\
MultiSeno & \texttt{RODADA-2/multi-seno-1\_0804\_19-06.csv} & 2 \\
SweptSine & \texttt{RODADA-4/swept-sine-1\_0819\_19-02.csv} & 4 \\
Degraus & \texttt{RODADA-2/seq-degraus-2\_0804\_19-38.csv} & 2 \\
Chirp & \texttt{RODADA-2/chirp-1\_0804\_19-19.csv} & 2 \\
\bottomrule
\end{tabular}
\end{table}

\begin{tcolorbox}[colback=blue!5!white,colframe=blue!75!black,title=Princípio de separação cross-session]
Nenhum arquivo de teste pertence à mesma rodada dos dados de treino dos Experimentos A--D (Rodada 5). Todos os testes são de sessões anteriores (Rodadas 2 e 4), garantindo avaliação em condições experimentais independentes.
\end{tcolorbox}

\section{Protocolo de Divisão dos Dados}

A divisão dos dados segue o princípio de \textbf{cross-session validation} recomendado na literatura de identificação de sistemas (Ljung, 1999; Schoukens \& Ljung, 2019):

\begin{figure}[hbt!]
\centering
\begin{tikzpicture}[
  box/.style={draw, rounded corners, align=center, minimum height=1cm, fill=gray!10},
  train/.style={box, fill=green!70!black, text=white},
  val/.style={box, fill=yellow!60!orange, text=black},
  test/.style={box, fill=orange!80!red, text=white},
  group/.style={draw, dashed, inner sep=10pt},
  >=stealth
]
\node[box] (r1) {Rodada 1\\31/07 -- 03/08};
\node[box, below=0.2cm of r1] (r2) {Rodada 2\\04/08};
\node[box, below=0.2cm of r2] (r3) {Rodada 3\\07/08};
\node[box, below=0.2cm of r3] (r4) {Rodada 4\\19/08};
\node[box, below=0.2cm of r4] (r5) {Rodada 5 (raiz)\\27/08};
\node[group, fit=(r1)(r2)(r3)(r4)(r5), label=above:Campanhas Experimentais] (rodadas) {};

\node[train, right=3cm of r3] (train) {TREINO\\(Exp A--D: Rodada 5)\\(Exp E: Rodada 3)};
\node[val, below=1cm of train] (val) {VALIDAÇÃO\\(Rodada 3)\\chirp-1\_16-32};
\node[test, below=1cm of val] (test) {TESTE\\(Rodadas 2 e 4)\\5 arquivos};
\node[group, fit=(train)(val)(test), label=above:Divisão dos Dados] (splits) {};

\draw[->] (r5.east) -- (train.west);
\draw[->] (r3.east) -- node[above, sloped, font=\scriptsize] {chirps excl. val} (train.west);
\draw[->] (r3.east) -- node[above, sloped, font=\scriptsize] {chirp-1\_16-32} (val.west);
\draw[->] (r2.east) -- (test.west);
\draw[->] (r4.east) -- (test.west);
\end{tikzpicture}
\caption{Protocolo de Divisão dos Dados}
\end{figure}

\subsection{Regras de Separação}

\begin{enumerate}
    \item \textbf{Nenhum arquivo aparece em mais de um split} (treino, validação ou teste).
    \item \textbf{O conjunto de teste é fixo e bloqueado}: jamais influencia decisões de treinamento ou seleção de modelo.
    \item \textbf{A validação guia exclusivamente o early stopping}, não é usada para ajuste de hiperparâmetros entre experimentos (todos usam os mesmos hiperparâmetros).
    \item \textbf{Separação temporal/por sessão}: treino e teste vêm de rodadas diferentes, minimizando viés por condições experimentais compartilhadas.
\end{enumerate}

\section{Pré-Processamento}

O pipeline de pré-processamento é idêntico para todos os conjuntos (treino, validação e teste):

\subsection{Etapas}

\begin{figure}[hbt!]
\centering
\begin{tikzpicture}[
  node distance=0.8cm,
  box/.style={draw, rectangle, align=center, minimum height=0.8cm, fill=blue!5},
  >=stealth
]
\node[box] (A) {CSV bruto\\($f_s = 100$ Hz)};
\node[box, below=of A] (B) {Filtrar referência > 0\\(remover repouso)};
\node[box, below=of B] (C) {Decimação $\times 2$\\$y$: filtro IIR anti-aliasing\\$u$: subsampling direto (ZOH)};
\node[box, below=of C] (D) {Recorte temporal\\$[200 : -200]$ amostras\\(remover transitórios)};
\node[box, below=of D] (E) {Normalização\\$\theta \to$ rad $\mid$ $u \to [-1, 1]$};
\node[box, below=of E] (F) {Derivada numérica\\Savitzky-Golay (janela=11, ordem=3)};
\node[box, below=of F] (G) {Estado final\\$x = [\theta \text{ (rad)}, d\theta/dt \text{ (rad/s)}]$};

\draw[->] (A) -- (B);
\draw[->] (B) -- (C);
\draw[->] (C) -- (D);
\draw[->] (D) -- (E);
\draw[->] (E) -- (F);
\draw[->] (F) -- (G);
\end{tikzpicture}
\caption{Etapas de Pré-Processamento}
\end{figure}

\subsection{Parâmetros}

\begin{table}[hbt!]
\centering
\caption{Parâmetros de Pré-Processamento}
\vspace{0.2cm}
\begin{tabular}{lp{3.5cm}p{7cm}}
\toprule
\textbf{Parâmetro} & \textbf{Valor} & \textbf{Justificativa} \\
\midrule
Fator de decimação & 2 & Reduz custo computacional mantendo banda útil ($f_{\text{Nyquist}}' = 25$ Hz $\gg$ banda do sistema) \\
Amostras cortadas (início) & 200 & Remove transitório inicial e artefatos de acionamento \\
Amostras cortadas (fim) & 200 & Remove artefatos de desligamento \\
Savitzky-Golay: janela & 11 amostras & Suavização moderada da derivada \\
Savitzky-Golay: ordem & 3 & Polinômio cúbico — adequado para capturar dinâmica suave \\
Normalização de $u$ & $u_{\text{norm}} = \text{clip}(u_{\%} / 100,\; -1,\; 1)$ & Escala unitária, saturação nos limites físicos \\
\bottomrule
\end{tabular}
\end{table}

\section{Protocolo de Treinamento}

\subsection{Variável Controlada}

Todos os seis experimentos utilizam:
\begin{itemize}
    \item \textbf{Mesmo modelo}: \texttt{PhysicsODE\_Asymmetric} (5 parâmetros)
    \item \textbf{Mesma inicialização}: \texttt{torch.manual\_seed(0)} e \texttt{np.random.seed(0)} resetados antes de cada experimento
    \item \textbf{Mesmos hiperparâmetros}: descritos na Tabela abaixo
\end{itemize}

A \textbf{única variável} entre experimentos é o \textbf{conjunto de dados de treino}.

\subsection{Hiperparâmetros (Fixos)}

\begin{table}[hbt!]
\centering
\caption{Hiperparâmetros (Fixos)}
\vspace{0.2cm}
\begin{tabular}{lcl}
\toprule
\textbf{Hiperparâmetro} & \textbf{Valor} & \textbf{Descrição} \\
\midrule
Épocas & 2000 & Número total de iterações de treinamento \\
Learning rate inicial & 0.015 & Taxa de aprendizado do Adam \\
Scheduler & CosineAnnealing & Decaimento suave do LR até 0 em $T_{\max} = 2000$ \\
$k_{\min}$ & 20 passos & Horizonte mínimo de predição (0.4 s) \\
$k_{\max}$ & 400 passos & Horizonte máximo de predição (8.0 s) \\
Estágio do currículo & 400 épocas & Intervalo entre dobras do horizonte \\
Batch size base & 1024 & Nº de janelas amostradas por época (total, todos os datasets) \\
Integrador & RK4 & Runge-Kutta 4ª ordem a passo fixo \\
Val check & cada 50 épocas & Frequência da avaliação de validação \\
\bottomrule
\end{tabular}
\end{table}

\subsection{Curriculum Learning}

O horizonte de predição $k$ cresce exponencialmente ao longo do treinamento, em estágios discretos:
%
\begin{equation}
k(\text{epoch}) = \min\left(k_{\min} \cdot 2^{\lfloor \text{epoch} / E_{\text{stage}} \rfloor},\; k_{\max}\right)
\end{equation}

\begin{table}[hbt!]
\centering
\caption{Curriculum Learning}
\vspace{0.2cm}
\begin{tabular}{ccccc}
\toprule
\textbf{Estágio} & \textbf{Épocas} & \textbf{$k$ (passos)} & \textbf{Horizonte temporal} & \textbf{Batch size efetivo} \\
\midrule
0 & 1--399 & 20 & 0.40 s & 1024 \\
1 & 400--799 & 40 & 0.80 s & 512 \\
2 & 800--1199 & 80 & 1.60 s & 256 \\
3 & 1200--1599 & 160 & 3.20 s & 128 \\
4 & 1600--2000 & 320 & 6.40 s & 64 \\
\bottomrule
\end{tabular}
\end{table}

\textbf{Motivação:} iniciar com horizontes curtos estabiliza os gradientes nas primeiras épocas (quando os parâmetros ainda estão longe dos valores corretos) e progressivamente força o modelo a capturar dinâmica de longo prazo (Chen et al., 2018; Turan et al., 2022).

O batch size é reduzido inversamente ao horizonte para manter a carga computacional aproximadamente constante.

\subsection{Função de Perda}

A perda é o erro quadrático médio normalizado pelo desvio padrão de cada estado, acumulado sobre todos os datasets de treino:
%
\begin{equation}
\mathcal{L} = \frac{1}{|\mathcal{D}|} \sum_{d \in \mathcal{D}} \frac{1}{B_d \cdot k} \sum_{i=1}^{B_d} \sum_{j=1}^{k} \left\| \frac{\hat{\mathbf{x}}_d^{(i)}(t_j) - \mathbf{x}_d^{(i)}(t_j)}{\boldsymbol{\sigma}_{\mathbf{x}}} \right\|^2
\end{equation}

onde:
\begin{itemize}
    \item $\mathcal{D}$ é o conjunto de datasets de treino do experimento
    \item $B_d$ é o número de janelas amostradas do dataset $d$
    \item $\hat{\mathbf{x}}$ é o estado predito pela integração da ODE
    \item $\mathbf{x}$ é o estado real (medido)
    \item $\boldsymbol{\sigma}_{\mathbf{x}} = [\sigma_\theta,\, \sigma_{\dot\theta}]$ é o desvio padrão de cada componente do estado calculado sobre todos os dados de treino
\end{itemize}

A normalização por $\boldsymbol{\sigma}_{\mathbf{x}}$ garante que a posição ($\theta$, tipicamente em dezenas de graus) e a velocidade ($\dot\theta$, tipicamente em rad/s) contribuam de forma equilibrada para a perda.

\subsection{Amostragem de Janelas}

Em cada época, para cada dataset $d$:
\begin{enumerate}
    \item Sorteia-se aleatoriamente $B_d = \max(1,\, \lfloor B_{\text{base}} / |\mathcal{D}| \rfloor)$ índices iniciais $i_0$
    \item Extrai-se a janela $\mathbf{x}_d[i_0 : i_0 + k]$ como alvo
    \item Integra-se a ODE a partir de $\mathbf{x}_d[i_0]$ por $k$ passos com a entrada $u_d$ interpolada
\end{enumerate}

\subsection{Early Stopping com Validação Free-Run}

A cada 50 épocas, avalia-se o modelo em \textbf{simulação livre completa} sobre o conjunto de validação (chirp da Rodada 3). Se o RMSE em graus for inferior ao melhor registrado, salva-se o checkpoint (\textit{state\_dict}). Ao final das 2000 épocas, o modelo é restaurado ao checkpoint com menor RMSE de validação.

\section{Definição dos Experimentos}

\begin{figure}[hbt!]
\centering
\begin{tikzpicture}[
  box/.style={draw, align=center, fill=gray!10, rounded corners, font=\scriptsize},
  group/.style={draw, dashed, inner sep=5pt},
  >=stealth,
  node distance=0.3cm and 0.5cm
]
\node[box] (A) {A: APRBS\\4 arquivos\\Rodada 5};
\node[box, right=of A] (B) {B: Multi-seno\\4 arquivos\\Rodada 5};
\node[box, right=of B] (C) {C: Swept-sine\\2 arquivos\\Rodada 5};
\node[box, below=of A] (D) {D: Seq-degraus\\4 arquivos\\Rodada 5};
\node[box, right=of D] (E) {E: Chirp\\3 arquivos\\Rodada 3};
\node[box, right=of E] (F) {F: Mix (todos)\\17 arquivos\\Rodadas 3+5};
\node[group, fit=(A)(B)(C)(D)(E)(F), label=above:6 Experimentos] (EXP) {};

\node[box, below=1cm of D] (VAL) {Val: chirp\\Rodada 3};
\node[box, right=1cm of VAL] (TEST) {Teste: 5 tipos\\Rodadas 2+4};

\node[box, fill=blue!10, below=2cm of EXP] (TRAIN) {train\_node()};
\node[box, fill=green!10, below=1cm of TRAIN] (EVAL) {avalia\_free\_run()};
\node[box, fill=orange!10, below=1cm of EVAL] (RES) {RMSE / R² / FIT\%\\por tipo de teste};

\draw[->] (EXP.south) -- node[right, font=\scriptsize] {mesmo modelo, seed, hparams} (TRAIN.north);
\draw[->] (VAL.east) -- node[above, font=\scriptsize] {early stopping} (TRAIN.west);
\draw[->] (TRAIN.south) -- node[right, font=\scriptsize] {modelo treinado} (EVAL.north);
\draw[->] (TEST.south) |- (EVAL.east);
\draw[->] (EVAL.south) -- (RES.north);
\end{tikzpicture}
\caption{Definição dos Experimentos}
\end{figure}

\begin{table}[hbt!]
\centering
\caption{Resumo dos Experimentos}
\vspace{0.2cm}
\begin{tabular}{clccl}
\toprule
\textbf{Exp} & \textbf{Treino} & \textbf{Nº Arq.} & \textbf{Rodada(s)} & \textbf{Propósito} \\
\midrule
A & Somente APRBS & 4 & 5 & Avaliar excitação binária pseudo-aleatória \\
B & Somente Multi-seno & 4 & 5 & Avaliar excitação multi-harmônica \\
C & Somente Swept-sine & 2 & 5 & Avaliar varredura contínua de frequência \\
D & Somente Seq-degraus & 4 & 5 & Avaliar excitação por transientes de degrau \\
E & Somente Chirp & 3 & 3 & Avaliar varredura logarítmica de frequência \\
F & Mix de A--E & 17 & 3 + 5 & \textbf{Controle} — diversidade vs. especialização \\
\bottomrule
\end{tabular}
\end{table}

\section{Avaliação e Métricas}

\subsection{Simulação Livre (\textit{Free-Run})}

A avaliação é feita em \textbf{simulação open-loop completa}: dada a condição inicial real $\mathbf{x}(t_0)$ e a sequência de entrada real $u(t)$, integra-se a ODE por toda a duração do dataset \textbf{sem qualquer correção ou reinicialização} do estado. Esta é a métrica definitiva para identificação de sistemas, pois testa a capacidade do modelo de reproduzir a dinâmica autonomamente (Ljung, 1999).

\subsection{Métricas}

Sejam $y_k = \theta_{\text{real}}(t_k)$ e $\hat{y}_k = \theta_{\text{pred}}(t_k)$ em graus, com $N$ amostras:

\subsubsection*{RMSE — Root Mean Square Error}
\begin{equation}
\text{RMSE} = \sqrt{\frac{1}{N} \sum_{k=1}^{N} (y_k - \hat{y}_k)^2} \quad [\text{graus}]
\end{equation}

\subsubsection*{R² — Coeficiente de Determinação}
\begin{equation}
R^2 = 1 - \frac{\sum_{k=1}^{N} (y_k - \hat{y}_k)^2}{\sum_{k=1}^{N} (y_k - \bar{y})^2}
\end{equation}

\subsubsection*{FIT\% — Normalized Root Mean Square Error}
\begin{equation}
\text{FIT} = \left(1 - \frac{\| \mathbf{y} - \hat{\mathbf{y}} \|}{\| \mathbf{y} - \bar{y}\,\mathbf{1} \|}\right) \times 100\%
\end{equation}

Equivalente à métrica utilizada pela função \texttt{compare} do MATLAB System Identification Toolbox.

\subsection{Matriz de Generalização}

Cada modelo treinado (linhas) é avaliado em cada tipo de teste (colunas), formando uma \textbf{matriz 6 × 5} de métricas:

\begin{table}[hbt!]
\centering
\caption{Matriz de Generalização}
\vspace{0.2cm}
\begin{tabular}{lccccc}
\toprule
 & \textbf{APRBS} & \textbf{MultiSeno} & \textbf{SweptSine} & \textbf{Degraus} & \textbf{Chirp} \\
\midrule
A\_APRBS & $\bullet$ & $\bullet$ & $\bullet$ & $\bullet$ & $\bullet$ \\
B\_MultiSeno & $\bullet$ & $\bullet$ & $\bullet$ & $\bullet$ & $\bullet$ \\
C\_SweptSine & $\bullet$ & $\bullet$ & $\bullet$ & $\bullet$ & $\bullet$ \\
D\_Degraus & $\bullet$ & $\bullet$ & $\bullet$ & $\bullet$ & $\bullet$ \\
E\_Chirp & $\bullet$ & $\bullet$ & $\bullet$ & $\bullet$ & $\bullet$ \\
F\_Mix & $\bullet$ & $\bullet$ & $\bullet$ & $\bullet$ & $\bullet$ \\
\bottomrule
\end{tabular}
\end{table}

Essa matriz permite identificar não apenas qual treino é melhor \textit{em média}, mas também padrões de especialização (e.g., treino X é superior no teste Y mas inferior em Z).

\subsection{Visualizações Geradas}

\begin{table}[hbt!]
\centering
\caption{Visualizações Geradas}
\vspace{0.2cm}
\begin{tabular}{lp{10cm}}
\toprule
\textbf{Arquivo} & \textbf{Conteúdo} \\
\midrule
\texttt{treino\_\{EXP\}.png} & Sinais de treino (ângulo e entrada vs. tempo) \\
\texttt{freerun\_\{EXP\}.png} & Sobreposição real vs. predito em todos os testes \\
\texttt{heatmap\_rmse.png} & Matriz 6$\times$5 colorida — RMSE \\
\texttt{heatmap\_r2.png} & Matriz 6$\times$5 colorida — R² \\
\texttt{heatmap\_fit.png} & Matriz 6$\times$5 colorida — FIT\% \\
\texttt{comparativo\_geral.png} & Barras de média geral (RMSE, R², FIT\%) por experimento \\
\texttt{resumo.json} & Todas as métricas em formato estruturado \\
\bottomrule
\end{tabular}
\end{table}

\section{Reprodutibilidade}

\begin{table}[hbt!]
\centering
\caption{Garantias de Reprodutibilidade}
\vspace{0.2cm}
\begin{tabular}{lp{11cm}}
\toprule
\textbf{Aspecto} & \textbf{Garantia} \\
\midrule
Seed determinística & \texttt{torch.manual\_seed(0)} + \texttt{np.random.seed(0)} resetados antes de cada experimento \\
Inicialização idêntica & Mesmo construtor \texttt{PhysicsODE\_Asymmetric()} com valores default \\
Dados versionados & Todos os CSVs hospedados no repositório GitHub com URLs fixas \\
Resultados salvos & Diretório \texttt{resultados\_v6\_\{timestamp\}/} com modelos \texttt{.pth}, plots \texttt{.png} e \texttt{resumo.json} \\
\bottomrule
\end{tabular}
\end{table}

\section{Limitações Conhecidas}

\begin{enumerate}
    \item \textbf{Desbalanceamento de dados}: o Experimento C (Swept-sine) possui apenas 2 arquivos enquanto os demais possuem 3--4. O Experimento F (Mix) possui 17 arquivos. Diferenças no volume de dados podem confundir o efeito do tipo de excitação com o efeito da quantidade de dados.
    \item \textbf{Seed única}: a ausência de múltiplas execuções com sementes diferentes impede a quantificação de intervalos de confiança. Recomenda-se rodar com 3--5 seeds e reportar média $\pm$ desvio padrão.
    \item \textbf{Validação com um único arquivo}: o early stopping depende de um único chirp (Rodada 3). Uma anomalia nesse arquivo específico enviesaria todos os experimentos igualmente.
    \item \textbf{Proximidade temporal entre val e treino no Exp E}: tanto o treino do Exp E quanto a validação provêm da Rodada 3 (mesma sessão, minutos de diferença), o que pode criar condições experimentais mais similares que para os demais experimentos.
\end{enumerate}

\section*{Referências}

\begin{enumerate}[label={[\arabic*]}]
    \item \textbf{Chen, R. T. Q., Rubanova, Y., Bettencourt, J., \& Duvenaud, D.} (2018). Neural Ordinary Differential Equations. \textit{Advances in Neural Information Processing Systems (NeurIPS)}, 31.
    \item \textbf{Ljung, L.} (1999). \textit{System Identification: Theory for the User} (2nd ed.). Prentice Hall.
    \item \textbf{Schoukens, J., \& Ljung, L.} (2019). Nonlinear System Identification: A User-Oriented Road Map. \textit{IEEE Control Systems Magazine}, 39(6), 28--99.
    \item \textbf{Kidger, P., Morrill, J., Foster, J., \& Lyons, T.} (2021). Neural Controlled Differential Equations for Irregular Time Series. \textit{Advances in Neural Information Processing Systems (NeurIPS)}, 34.
    \item \textbf{Turan, C., Shaker, H. R., \& Shardt, Y. A. W.} (2022). Physics-informed neural ODEs for learning process dynamics. \textit{Journal of Process Control}, 120, 1--12.
    \item \textbf{Rackauckas, C., Ma, Y., Martensen, J., et al.} (2020). Universal Differential Equations for Scientific Machine Learning. \textit{arXiv preprint arXiv:2001.04385}.
\end{enumerate}

\end{document}
